% ID: SEL-D0063-OA-E03
% ESTADO: cotejada. Segunda pasada visual y borrador compilado revisados.
% Fuente: Selectividad/sel_2010_matematicas_aplicadas/mat_apli/mat_apli_junio/mat_apli_junio_rb/EXAMEN-ANDALUCÍA-6.pdf
% SHA-256: 0de2f7d077a3b5d5c0dbca886cdb363082356411843369c3bce54ff7cda08abc
% Página: 1; opción A; ejercicio original 3.

En el experimento aleatorio consistente en lanzar un dado equilibrado con las caras numeradas del 1 al 6 y observar el resultado se consideran los siguientes sucesos:

$A$: ``obtener un número mayor que 4'', $B$: ``obtener un número par''.
\begin{enumerate}[label=\alph*)]
\item \textbf{(1 punto)} Escriba los elementos de cada uno de los siguientes sucesos:
\[
A;\quad B;\quad A^C\cup B;\quad A\cap B^C;\quad(A\cap B)^C.
\]
\item \textbf{(1.5 puntos)} Calcule las probabilidades $P\bigl(A^C\cap B^C\bigr)$ y $P\bigl(A^C\cup B^C\bigr)$.
\end{enumerate}
